\documentclass[11pt,a4paper]{article}

% ---------- fontes e língua ----------
\usepackage[T1]{fontenc}
\usepackage[utf8]{inputenc}
\usepackage{lmodern}
\usepackage{tgpagella}
\usepackage{tgheros}
\usepackage[portuguese]{babel}
\usepackage{microtype}

% ---------- página e cores ----------
\usepackage[a4paper,top=2.6cm,bottom=2.4cm,left=2.4cm,right=2.4cm]{geometry}
\usepackage[table]{xcolor}
\definecolor{azul}{RGB}{28,62,112}
\definecolor{azulclaro}{RGB}{226,234,246}
\definecolor{cinza}{RGB}{95,102,115}

% ---------- tabelas, listas, caixas ----------
\usepackage{booktabs}
\usepackage{array}
\usepackage{tabularx}
\usepackage{enumitem}
\usepackage{amssymb}
\usepackage{titlesec}
\usepackage{fancyhdr}
\usepackage{lastpage}
\usepackage{tcolorbox}

\setlength{\parindent}{0pt}
\setlength{\parskip}{6pt}

\setlist{nosep, leftmargin=*, itemsep=3pt}
\setlist[itemize]{label=\textcolor{azul}{\textbullet}}
\setlist[enumerate]{label=\textcolor{azul}{\sffamily\bfseries\arabic*.}}

\titleformat{\section}{\sffamily\Large\bfseries\color{azul}}{\thesection}{0.7em}{}[{\color{azul}\titlerule[0.6pt]}]
\titleformat{\subsection}{\sffamily\large\bfseries\color{azul}}{\thesubsection}{0.6em}{}
\titlespacing*{\section}{0pt}{1.6em}{0.8em}
\titlespacing*{\subsection}{0pt}{1.1em}{0.3em}

\pagestyle{fancy}
\setlength{\headheight}{15pt}
\fancyhf{}
\renewcommand{\headrule}{{\color{azul}\hrule width\headwidth height\headrulewidth}}
\renewcommand{\headrulewidth}{0.4pt}
\fancyhead[L]{\sffamily\small\color{cinza} Universidade da Beira Interior}
\fancyhead[R]{\sffamily\small\color{cinza} Engenharia de Software \textperiodcentered{} 2026/27}
\fancyfoot[L]{\sffamily\small\color{cinza} Contrato de Equipa}
\fancyfoot[R]{\sffamily\small\color{cinza} Página \thepage\ de \pageref{LastPage}}

\fancypagestyle{primeira}{%
  \fancyhf{}%
  \renewcommand{\headrulewidth}{0pt}%
  \fancyfoot[L]{\sffamily\small\color{cinza} Contrato de Equipa}%
  \fancyfoot[R]{\sffamily\small\color{cinza} Página \thepage\ de \pageref{LastPage}}%
}

\newcolumntype{Y}{>{\raggedright\arraybackslash}X}
\newcolumntype{L}[1]{>{\raggedright\arraybackslash}p{#1}}
\newcolumntype{N}[1]{>{\raggedright\arraybackslash\sffamily\bfseries}p{#1}}
\newcommand{\th}[1]{\textcolor{white}{\sffamily\bfseries #1}}

\tcbset{
  caixa/.style={colback=white, colframe=azul, coltitle=white, colbacktitle=azul,
    fonttitle=\sffamily\bfseries, arc=3pt, boxrule=0.6pt,
    left=9pt, right=9pt, top=6pt, bottom=6pt}
}

\begin{document}
\thispagestyle{primeira}

% ---------- cabeçalho do documento ----------
\begin{tcolorbox}[colback=azul, colframe=azul, coltext=white, arc=3pt, boxrule=0pt,
    left=16pt, right=16pt, top=14pt, bottom=14pt]
  \sffamily
  {\footnotesize\bfseries\MakeUppercase{Universidade da Beira Interior}}\\[2pt]
  {\footnotesize Faculdade de Engenharia \textperiodcentered{} Departamento de Informática}\\[16pt]
  {\fontsize{28}{32}\selectfont\bfseries Contrato de Equipa}\\[4pt]
  {\large Team Contract \textperiodcentered{} Engenharia de Software 2026/27}
\end{tcolorbox}

Este contrato regista o que combinámos para trabalhar juntos durante o projeto de Engenharia de Software, que decorre em cinco sprints a partir de 21 de setembro de 2026. Serve para evitar mal-entendidos, para cada um saber o que se espera de si e para termos um sítio a que recorrer quando houver dúvidas. Está também no quadro Jira, na secção \textit{Team Contract}.

\begin{tcolorbox}[caixa, title={Em resumo}, colback=azulclaro!60]
  \begin{itemize}
    \item Daily Scrum todos os dias úteis (15 minutos), Planning no primeiro dia do sprint, Review e Retrospetiva no fim.
    \item Em cada sprint há um Product Owner e um Scrum Master, e todos passamos pelos dois papéis.
    \item Decidimos por consenso. Se não der, votamos.
    \item Nada fica feito sem revisão, testes e aceitação do Product Owner.
    \item Tudo passa por \textit{pull request} na \texttt{dev}; a \texttt{main} só recebe código no fim do sprint.
    \item A IA ajuda, mas cada um responde pelo que entrega, e o uso fica registado.
  \end{itemize}
\end{tcolorbox}

\section*{A equipa}

\arrayrulecolor{azul!35}
\renewcommand{\arraystretch}{2.2}
\begin{tabularx}{\textwidth}{Y L{5.2cm}}
  \rowcolor{azul} \th{Nome} & \th{N.º de aluno} \\
  \rowcolor{azulclaro!55} Felipe Fogaça   & \\ \hline
  \rowcolor{white}        João Vitor      & \\ \hline
  \rowcolor{azulclaro!55} Marcus Bernardo & \\ \hline
  \rowcolor{white}        Rui Machado     & \\ \hline
  \rowcolor{azulclaro!55} Tiago Pereira   & \\ \hline
\end{tabularx}
\renewcommand{\arraystretch}{1}
\arrayrulecolor{black}

%--------------------------------------------------------------------
\section{Reuniões}

\subsection{As reuniões que fazemos}

Seguimos o Scrum e fazemos as reuniões habituais. As durações são máximos: se acabarmos mais cedo, melhor.

\small
\renewcommand{\arraystretch}{1.45}
\begin{tabularx}{\textwidth}{N{2.5cm} L{3.2cm} L{1.7cm} Y}
  \rowcolor{azul} \th{Reunião} & \th{Quando} & \th{Duração} & \th{O que fazemos} \\
  \rowcolor{azulclaro!55}
  Sprint Planning & No primeiro dia de cada sprint & até 1h &
    O Product Owner traz o backlog já priorizado. Definimos o objetivo do sprint, escolhemos os itens que cabem, dividimos em tarefas, estimamos em story points e cada tarefa fica com um responsável. \\
  \rowcolor{white}
  Daily Scrum & Todos os dias úteis, a uma hora que dê para todos & 15 min &
    Cada um diz o que fez desde o último dia, o que vai fazer hoje e o que o está a bloquear. \\
  \rowcolor{azulclaro!55}
  Refinement & Uma vez por semana, a meio da semana & 30 min &
    O Product Owner e a equipa afinam os próximos itens do backlog: critérios de aceitação, estimativas e dúvidas. \\
  \rowcolor{white}
  Sprint Review & No fim do sprint, antes da apresentação oral & até 1h &
    Mostramos o que ficou pronto e vemos se o objetivo do sprint foi cumprido. O Product Owner vê se aceita ou não os itens. \\
  \rowcolor{azulclaro!55}
  Retrospetiva & Logo a seguir à Review & 45 min &
    Falamos do que correu bem, do que correu mal e escolhemos duas ou três melhorias concretas para o sprint seguinte, cada uma com um membro responsável para melhorar, rotativo. \\
  \bottomrule
\end{tabularx}
\renewcommand{\arraystretch}{1}
\normalsize

No Sprint 5, que dura cinco semanas, fazemos ainda um ponto de situação de 30 minutos a meio, para perceber se estamos dentro do plano.

\subsection{Como registamos no Jira}

O enunciado pede provas de que cumprimos o Scrum, por isso tudo fica escrito na secção \textit{Daily SCRUM} do quadro:

\begin{itemize}
  \item \textbf{Daily:} cada um escreve as suas três linhas (feito, a fazer, bloqueios), com a data. O Scrum Master confirma, até ao fim do dia, que está tudo escrito.
  \item \textbf{Planning, Review e Retrospetiva:} o Scrum Master escreve uma ata curta (data, quem esteve, decisões e ações com responsável).
\end{itemize}

\subsection{Presenças}

\begin{itemize}
  \item Chegamos a horas, com 5 minutos de tolerância.
  \item Se alguém não puder estar, avisa no grupo com tempo considerável de antecedência.
  \item Quem falta lê o registo no Jira e escreve o seu contributo no mesmo dia.
  \item Se não der para estarmos juntos no daily, fazemo-lo por escrito no Jira, ou por mensagem no grupo de WhatsApp.
  \item Durante as reuniões falamos um de cada vez, deixamos o telemóvel de lado e o que fugir ao tema fica anotado para depois.
  \item As aulas presenciais também contam: convém estar todos, como o enunciado pede.
\end{itemize}

%--------------------------------------------------------------------
\section{Responsabilidades}

\subsection{Da equipa}

\begin{itemize}
  \item Entregar, no fim de cada sprint, um incremento que funcione e apresentá-lo oralmente na última semana do sprint.
  \item Manter o quadro Jira atualizado (\textit{Product Backlog}, \textit{To Do}, \textit{In Progress}, \textit{Done}, \textit{Artifacts}, \textit{Daily SCRUM} e \textit{Team Contract}), com acesso para o professor.
  \item Decidir por consenso. Se não chegarmos lá, votamos e vale a maioria, a decisão fica registada no Jira.
  \item Repartir o trabalho de forma justa. O projeto tem \textit{back-end} e \textit{front-end} e a ideia é que toda a gente passe pelos dois, para nos podermos cobrir uns aos outros.
  \item Garantir que todos percebem o que está a ser feito: qualquer um de nós tem de saber explicar o projeto nas apresentações.
\end{itemize}

\subsection{De cada elemento}

\begin{itemize}
  \item Ir às reuniões e às aulas presenciais.
  \item Fazer as suas tarefas dentro do prazo e mover o cartão no Jira no momento em que o estado muda.
  \item Avisar cedo quando está bloqueado ou vai atrasar-se.
  \item Rever os \textit{pull requests}.
  \item Cumprir a Definição de Concluído e as regras deste contrato.
\end{itemize}

\subsection{Product Owner e Scrum Master (Team Captain)}
Em cada sprint há um Product Owner e um Scrum Master. Ambos continuam a desenvolver como os outros; no Planning ajustamos a carga de tarefas para lhes sobrar tempo para o papel.

\begin{tcolorbox}[caixa, title={Product Owner}]
  \begin{itemize}
    \item Gere e prioriza o Product Backlog, com o que traz mais valor no topo da lista.
    \item Explica os itens à equipa e ajuda a escrever os critérios de aceitação.
    \item Está disponível durante o sprint para tirar dúvidas.
    \item Na Review decide se o objetivo do sprint foi cumprido e aceita ou recusa os itens.
    \item Não mexe no sprint backlog a meio do sprint sem falar com a equipa.
  \end{itemize}
\end{tcolorbox}

\begin{tcolorbox}[caixa, title={Scrum Master}]
  \begin{itemize}
    \item Garante que as reuniões acontecem, e que ficam registadas no Jira.
    \item Ajuda a equipa a usar bem o Scrum.
    \item Ajuda a remover os obstáculos que apareçam e leva aos docentes o que a equipa não conseguir resolver sozinha.
    \item Faz de mediador quando há desacordos.
    \item Coordena o resumo para a apresentação oral do sprint.
  \end{itemize}
\end{tcolorbox}

\subsection{Rotação nos cinco sprints}

Cada elemento é Product Owner uma vez e Scrum Master uma vez, e ninguém tem os dois papéis no mesmo sprint.

\renewcommand{\arraystretch}{1.5}
\begin{tabularx}{\textwidth}{N{1.6cm} L{5.2cm} Y Y}
  \rowcolor{azul} \th{Sprint} & \th{Datas} & \th{Product Owner} & \th{Scrum Master} \\
  \rowcolor{azulclaro!55} 1 & 21 de setembro a 14 de outubro & Tiago Pereira  & Felipe Fogaça \\
  \rowcolor{white}        2 & 12 a 30 de outubro             & Felipe Fogaça   & João Vitor \\
  \rowcolor{azulclaro!55} 3 & 2 a 20 de novembro             & João Vitor      & Marcus Bernardo \\
  \rowcolor{white}        4 & 23 de novembro a 11 de dezembro & Marcus Bernardo & Rui Machado \\
  \rowcolor{azulclaro!55} 5 & 14 de dezembro a 15 de janeiro & Rui Machado     & Tiago Pereira \\
  \bottomrule
\end{tabularx}
\renewcommand{\arraystretch}{1}

Só trocamos de papéis se todos estiverem de acordo, e a troca fica registada no Jira antes do Sprint Planning.

%--------------------------------------------------------------------
\section{Conflitos}

\subsection{Como os resolvemos}

\begin{enumerate}
  \item \textbf{Falar diretamente.} As pessoas envolvidas conversam com calma e o mais cedo possível.
  \item \textbf{Levar à equipa.} Se não se resolver, trazemos o assunto para uma reunião curta ou para a Retrospetiva, com o Scrum Master como mediador. Se for o Scrum Master a estar envolvido, quem medeia é o Product Owner.
  \item \textbf{Decidir.} Procuramos consenso. Se não houver, votamos e a decisão fica no Jira. Todos cumprem o que ficou decidido, mesmo quem votou contra.
  \item \textbf{Pedir ajuda.} Se o problema continuar ou se repetir, falamos com os docentes.
\end{enumerate}

\subsection{Quando alguém não cumpre}

Falhas que contam: faltar a reuniões ou aulas sem justificação, não entregar tarefas sem avisar, não atualizar o Jira ou não respeitar as regras do GitHub e a Definição de Concluído.

\begin{enumerate}
  \item \textbf{Primeira vez:} conversa informal com o Scrum Master para recordar o que combinámos. Fica uma nota curta no Jira.
  \item \textbf{Segunda vez:} aviso escrito no Jira com um novo prazo. A tarefa pode ser passada a outro colega e a pessoa compensa com trabalho no sprint seguinte.
  \item \textbf{Terceira vez:} informamos os docentes. A contribuição de cada um está documentada no Jira e no GitHub, por isso a avaliação individual fica justa para quem cumpriu.
\end{enumerate}


%--------------------------------------------------------------------
\section{Outras regras}

\subsection{Comunicação}

\begin{itemize}
  \item Usamos um grupo de WhatsApp para o dia a dia e avisos rápidos.
  \item Tudo o que seja decisão, tarefa ou requisito passa para o Jira. O que ficar só na conversa pode perder-se.
\end{itemize}

\subsection{Prazos}

\begin{itemize}
  \item As datas dos sprints são as que os docentes definiram e a entrega final é no último dia do Sprint 5.
  \item Damos o trabalho por fechado um dia útil antes do fim do sprint, para haver tempo para a Review, a Retrospetiva e a apresentação oral.
  \item Uma tarefa só entra em \textit{In Progress} depois de ter responsável e estimativa, e ninguém tem mais de duas tarefas abertas ao mesmo tempo.
  \item Se um item não ficar pronto, não prolongamos o sprint: volta ao Product Backlog e o Product Owner volta a priorizá-lo.
\end{itemize}

\subsection{Definição de Concluído}

Um item só passa a \textit{Done} quando:

\begin{tcolorbox}[caixa, title={Definição de Concluído}, colback=azulclaro!40]
  \begin{itemize}[label=\textcolor{azul}{$\square$}]
    \item cumpre os critérios de aceitação combinados com o Product Owner;
    \item o código foi revisto por pelo menos outro elemento da equipa (\textit{pull request} aprovado);
    \item tem testes (unitários e, quando fizer sentido, de integração) e todos passam;
    \item está integrado na branch \texttt{dev} sem quebrar o build nem a integração contínua;
    \item a documentação necessária está atualizada e os artefactos estão ligados no Jira;
    \item o Product Owner aceitou o item;
    \item o cartão no Jira está em \textit{Done}, com a ligação para o \textit{pull request}.
  \end{itemize}
\end{tcolorbox}

\subsection{Uso do GitHub}

Temos um repositório da equipa, ligado no Jira na secção \textit{Artifacts}, e trabalhamos com três tipos de branch:

\small
\renewcommand{\arraystretch}{1.5}
\begin{tabularx}{\textwidth}{N{4.9cm} Y}
  \rowcolor{azul} \th{Branch} & \th{Para que serve} \\
  \rowcolor{azulclaro!55} \texttt{main} & Versões estáveis. Só recebe código no fim do sprint, depois da Review. \\
  \rowcolor{white} \texttt{dev} & Onde integramos o trabalho de todos. \\
  \rowcolor{azulclaro!55} \texttt{feature/PROJ-12-login} & Uma por tarefa, criada a partir da \texttt{dev}, com o código do cartão do Jira no nome. \\
  \bottomrule
\end{tabularx}
\renewcommand{\arraystretch}{1}
\normalsize

\begin{itemize}
  \item Fazemos commits pequenos e frequentes, com mensagens curtas que digam o que mudou.
  \item Passwords, chaves de API e dados pessoais não entram no repositório.
  \item Antes de começar a trabalhar fazemos \textit{pull} da \texttt{dev}.
  \item Se a integração contínua falhar, corrigir isso passa à frente do resto.
\end{itemize}

\subsection{Uso responsável de IA}

O projeto incentiva o uso de ferramentas de IA e nós vamos usá-las, mas com regras:

\begin{itemize}
  \item A equipa escolhe pelo menos uma ferramenta de IA para apoiar a programação e regista no Jira qual é.
  \item A IA ajuda, mas não decide por nós. Só entra no projeto o que percebemos e conseguimos explicar nas apresentações: cada um responde pelo que entrega.
  \item O que for gerado por IA é revisto e testado como qualquer outro código ou texto, e passa pela mesma Definição de Concluído.
  \item Não pomos passwords, chaves, dados pessoais nem informação que não possamos partilhar em ferramentas de IA.
  \item Verificamos sempre a informação e o código que a IA sugere, incluindo licenças e fontes.
  \item Registamos onde e para quê usámos IA no ficheiro \textit{Registo de Uso de IA}, que está ligado no Jira (\textit{Artifacts}), e assumimos quando algo foi gerado. A base deste contrato foi feita com apoio de IA e está declarada nesse registo.
\end{itemize}

%--------------------------------------------------------------------
\section{Alterações e confirmação}

Revemos este contrato nas retrospetivas. Qualquer alteração precisa da concordância de todos e fica registada no Jira. Ao assinar, cada elemento da equipa confirma que leu o contrato e que está de acordo.

\vspace{0.4em}
\begin{tabularx}{\textwidth}{Y@{\hspace{1.4cm}}Y}
  \rule{0pt}{2.6em}\hrulefill & \rule{0pt}{2.6em}\hrulefill \\
  \small Felipe Fogaça & \small João Vitor \\
  \rule{0pt}{2.6em}\hrulefill & \rule{0pt}{2.6em}\hrulefill \\
  \small Marcus Bernardo & \small Rui Machado \\
  \rule{0pt}{2.6em}\hrulefill & \\
  \small Tiago Pereira & \\
\end{tabularx}

\vspace{1.5em}
Covilhã, \underline{\hspace{1cm}} de outubro de 2026

\end{document}
